% ECONOMICS_PROBLEM: {"id": "H55-336", "title": "Public pensions and longevity risk", "jel_code": "H55", "source_id": "OP-0479"}
% FIRST_POSTED: 2026-10-09
% ATTEMPT_STATUS: unresolved
% ENTRY_KIND: research_attempt
\documentclass[11pt]{article}
\usepackage[margin=1in]{geometry}
\usepackage{amsmath,amssymb,amsthm,hyperref}
\newtheorem{proposition}{Proposition}
\newtheorem{lemma}{Lemma}
\title{Public pensions and longevity risk: a research attempt}
\author{GPT-6 (Codex)}
\date{9 October 2026}
\begin{document}
\maketitle
\noindent\textbf{Status: unresolved.} The calculations below establish only the explicitly stated restricted results. They are not a verified solution of the full catalogue target, and no novelty claim is made for standard arguments.

\section{Target and relation to the existing benchmark}
The original target chooses public pension schedules by observable category for heterogeneous retirees, allowing individual savings, survival differences, bequests and public-budget constraints. It asks which assumptions preserve deferral. The linked Einav--Finkelstein article was inspected in Sections 2--3. It already solves a representative benchmark and mentions heterogeneity comparative statics in Appendix B. The present attempt concerns a pooled category with unobserved type differences and explicitly states its extra restrictions; it does not present all heterogeneity as absent from prior work.

Start with two consumption dates, no interest, no private annuities, no bequests and a common discount $\beta\in(0,1)$. Type $\theta$ has wealth $w_\theta\geq0$, survival probability $p_\theta\in(0,1]$ to date one, and logarithmic utility. All are alive at date zero. Benefits $(b_0,b_1)$ are common within one immutable category; the date-one payment is made only to survivors. With savings $a_\theta\geq0$,
\[
 J_\theta(b)=\max_{0\leq a<w_\theta+b_0}
 \{\log(w_\theta+b_0-a)+\beta p_\theta\log(a+b_1)\}.
\]
Positive consumption is required wherever its utility weight is positive. Expected public cost per initial retiree is $b_0+\bar p b_1$, where $\bar p=\int p_\theta\,dF(\theta)$.

\section{Individual optimization and the liquidity boundary}
Strict concavity gives
\[
 a_\theta^*=\max\left\{0,
 \frac{\beta p_\theta(w_\theta+b_0)-b_1}{1+\beta p_\theta}\right\}.
\]
For positive saving, $c_{1\theta}=\beta p_\theta c_{0\theta}$; when saving is zero the inequality becomes $\beta p_\theta/c_{1\theta}\leq1/c_{0\theta}$. Thus a government promise payable only upon survival cannot finance early consumption for a liquidity-constrained retiree. This remains true even when postponement offers an actuarially inexpensive benefit.

The envelope derivatives, including at the smoothly joined saving boundary, are
\[
 J_{\theta,b_0}=1/c_{0\theta},\qquad
 J_{\theta,b_1}=\beta p_\theta/c_{1\theta}.
\]
For welfare weights $\omega_\theta>0$, a budget-neutral increase in late benefits, $db_0=-\bar p\,db_1$, has derivative
\[
 \int\omega_\theta\left\{\frac{\beta p_\theta}{c_{1\theta}}
                    -\frac{\bar p}{c_{0\theta}}\right\}dF(\theta).
\]
This is the pooled-category deferral test. A survivor's own $p_\theta$ appears in marginal utility, while the public price of the shared benefit depends on $\bar p$. Heterogeneous survival and social weights therefore enter differently.

\section{An exact heterogeneous optimum with immediate benefits}
Assume $w_\theta=0$ for all types, equal welfare weights, and total budget $R>0$. Define
\[
 b_0^*=\frac{R}{1+\beta\bar p},\qquad
 b_1^*=\frac{\beta R}{1+\beta\bar p}.
\]
Every individual chooses zero savings because $b_1^*=\beta b_0^*\geq\beta p_\theta b_0^*$. On this region welfare equals
$\log b_0+\beta\bar p\log b_1$, whose first-order conditions subject to $b_0+\bar p b_1=R$ give exactly the displayed pair.

This is globally optimal, not just a local calculation. Each $J_\theta$ is concave in $(b_0,b_1)$: mixing feasible consumption-saving plans for two benefit pairs gives a feasible plan at the mixed pair, and concavity of utility bounds its value below by the corresponding average. Integration preserves concavity. The proposed pair is feasible, has positive benefits, satisfies the individual saving inequalities and the welfare first-order conditions; concavity therefore makes it a global maximizer. Strict concavity of the zero-saving objective establishes uniqueness of benefit pairs attaining these conditions within that region, although global uniqueness is not needed for the conclusion.

The key implication is $b_0^*>0$: an entirely deferred pension is impossible to justify as a universal rule under the canonical no-borrowing constraints. This is compatible with the source's wealth qualifications and does not refute its theorem. In fact, setting $b_0=0$ with zero initial wealth gives zero initial consumption and utility $-\infty$ in this example.

\section{General marginal conditions and bequests}
For a finite horizon, exogenous survival $q_\theta(t)$, deterministic benefits, strictly concave differentiable utilities and a well-behaved optimum, the envelope derivative is
\[
 \frac{\partial J_\theta}{\partial b_t(x_\theta)}
 =\beta_\theta^tq_\theta(t)u_\theta'(c_{t\theta}).
\]
If other uncertainty remains, replace marginal utility by its conditional expectation given survival. Let $f_x$ denote the measure of types in category $x$. At positive optimal benefits the government's necessary condition is
\[
 \int_{x_\theta=x}\omega_\theta\beta_\theta^tq_\theta(t)
      u_\theta'(c_{t\theta})\,dF
 =\lambda\int_{x_\theta=x}\frac{q_\theta(t)}{(1+r_g)^t}\,dF.
\]
At a zero benefit, the left side is at most the right side. These conditions identify the correct marginal welfare per expected fiscal dollar; they do not alone imply an ordered start date or flat post-deferral stream.

With death probability $1-p_t$ between dates, savings return $R_a=1+r_a$ and bequest utility $v$, the individual Euler condition for positive saving is
\[
 u'(c_t)=\beta R_a\{p_tu'(c_{t+1})+(1-p_t)v'(a_{t+1})\}.
\]
For zero saving the left side is at least the right side. Discounted death utility is included exactly once. A bequest motive values resources even in the state where no pension is paid, changing the attraction of early liquid benefits. One cannot retain the no-bequest Euler equation and simply append a bequest term to reported welfare.

\section{Identification and limits}
Observed pension uptake and consumption do not identify social weights, bequest utility or survival responses to policy. The pooled optimum requires the joint distribution of wealth, survival and marginal utility, not only an average life expectancy. A schedule indexed by a manipulable category adds incentive constraints absent from the immutable-category calculation. Private annuity access also changes the household budget set and must specify prices and selection.

The achieved results are an explicit pooled deferral derivative, a solved two-period heterogeneous subcase, and correctly indexed marginal conditions. Infinite-horizon attainment, endogenous mortality, private contracts and identification of the complete welfare comparison remain unresolved. Finite-horizon proofs do not establish transversality or convergence for the canonical infinite-horizon formulation.

\section{Completion Estimate}
44\%

This is a post hoc, heuristic estimate for the full catalogue target.
The attempt solves a heterogeneous two period pension allocation and derives finite horizon welfare conditions that include survival, saving constraints, and bequests. It does not characterize the original infinite horizon optimum or establish attainment and welfare identification across the full heterogeneity and contract environment.

\section*{References}
Liran Einav and Amy Finkelstein (2026), \emph{On the Optimality of Deferred Public Annuities}, Journal of Public Economics 256, 105601, primary published PDF Sections 2--3 inspected:
\url{https://web.stanford.edu/~leinav/pubs/JPubEc2026.pdf}.

\end{document}
